% ===========================================================================
% Self-contained SM section, placed after "Robustness" (supp:robustness) and before "Supplementary videos".
% The tables tab_sm_*.tex are written by scripts/em/f9_sm_tablas.py (make em-tables) from results_em/.
\section{Adult recovery, mate search and colony size}
\label{supp:recovery}\label{sec:recovery}

This section gives, from archived runs, the numbers on which the
conditions of the mechanism rest (main text, Sections~3.2 and~4.2): the extinction probability as a function of adult
recovery, its dependence on colony size, the effect of removing non-breeders, of background mortality and of
protecting the juveniles, the test of a mate-finding Allee effect, the scaling of the extinction
time, the cohort-by-cohort account of the mechanism in the well-mixed limit, and the held-out patterns of the
reversible branch. Every count is reproduced by
\texttt{scripts/em/f9\_sm\_tablas.py} in the code repository.

\paragraph{Design}
The adult recovery factor $\eadult\in[0,1]$ multiplies the gain of an individual outside the window: an adult gains
$\eadult\,r[(1-\lambda)+\lambda\bar s_V]$ per week, so that with $\lambda<1$ a non-breeding adult recovers on its own
at rate $\eadult r(1-\lambda)$, and with $\lambda=1$ only from competent neighbours. \Cone\ is $\eadult=0$, $\lambda=1$.
We ran $\eadult\in\{0,0.01,0.02,0.05,0.1,0.2,0.3,0.4,0.5,0.6,0.8,1\}$ and $\lambda\in\{0,0.5,0.8,0.9,0.95,1\}$, with
all other rules of \Cone, at $L=15$ (20 replicates), $L=30$ (20) and $L=60$ (10; nine values of $\eadult$ and a
subset of $\lambda$), in space (with retreat sites) and well mixed (without retreat sites; father drawn in the
female's cell, in her $3\times3$ block or in the whole arena): $12\,870$ colonies in all, with the founders and the
seeding density of \Cone. The response is the probability of extinction by week $1\,000$, with $95\%$ Wilson
intervals \citep{Wilson1927}; for an ensemble in which every colony goes extinct the interval is $[0.84,1]$ with 20
replicates and $[0.72,1]$ with 10. The threshold $\eadult^c(\lambda,L)$ is the midpoint of a binomial
maximum-likelihood logistic fit of the extinction probability against $\eadult$.

\paragraph{Small adult recovery does not prevent the collapse}
With $\eadult\le0.02$ every colony went extinct: $2\,880$ of $2\,880$ at $L=15$ and $30$, in space and well mixed,
under the three mating rules and for every $\lambda$ (Table~\ref{tab:sm-pext}); with $\eadult=0$, every one of the
$1\,110$ colonies at $L=15$--$60$. The collapse is therefore not an artefact of the edge $\eadult=0$. It is lost
when adults recover fast enough on their own (Table~\ref{tab:sm-epsc}). At the size of the experiment ($L=30$)
the threshold is $0.085$ in space with individual recovery ($\lambda=0$) and $0.73$ in space with purely social
recovery ($\lambda=1$); well mixed it is $0.70$ for $\lambda=0$ and exceeds $1$ for $\lambda\ge0.5$ (at $\lambda=1$
every colony goes extinct even with $\eadult=1$): without competent demonstrators an adult whose recovery is social
never recovers. The relevant condition is thus that
adults do not recover \emph{asocially}, that is, that $\eadult(1-\lambda)$ is small, and in space that the pockets of
competent animals that survive the crowded phase are too few to teach the adults around them.

\begin{table}[ht]
\centering
\caption{Probability that a colony of \Cone\ with adult recovery factor $\eadult$ is extinct by week $1\,000$, for
purely social ($\lambda=1$) and purely individual ($\lambda=0$) learning and recovery, in space (with retreat sites)
and well mixed (no retreat sites, mating in the cell), at $L=15$, $30$ and $60$. $n$: replicates per cell of the
table. --: not simulated. 95\% Wilson intervals: $[0.84,1]$ for $20/20$, $[0.72,1]$ for $10/10$, at most $\pm0.22$
otherwise. Source: \texttt{results\_em/adult\_recovery/summary.csv}.}
\label{tab:sm-pext}
\scriptsize
\setlength{\tabcolsep}{2.6pt}
\begin{tabular}{@{}lcrr*{12}{c}@{}}
\toprule
& & & & \multicolumn{12}{c}{$\eadult$} \\
\cmidrule(l){5-16}
Structure & $\lambda$ & $L$ & $n$ & 0 & .01 & .02 & .05 & .1 & .2 & .3 & .4 & .5 & .6 & .8 & 1 \\
\midrule
spatial & 1 & 15 & 20 & 1 & 1 & 1 & 1 & 1 & 1 & 1 & .95 & .95 & .95 & .60 & .55 \\
spatial & 1 & 30 & 20 & 1 & 1 & 1 & 1 & 1 & 1 & .95 & .90 & .85 & .70 & .50 & .05 \\
spatial & 1 & 60 & 10 & 1 & -- & -- & 1 & 1 & .90 & 1 & .80 & -- & .40 & 0 & 0 \\
\addlinespace
spatial & 0 & 15 & 20 & 1 & 1 & 1 & 1 & .80 & .05 & 0 & 0 & 0 & 0 & 0 & 0 \\
spatial & 0 & 30 & 20 & 1 & 1 & 1 & .85 & .30 & 0 & 0 & 0 & 0 & 0 & 0 & 0 \\
spatial & 0 & 60 & 10 & 1 & -- & -- & .50 & 0 & 0 & 0 & 0 & -- & 0 & 0 & 0 \\
\addlinespace
well mixed & 1 & 15 & 20 & 1 & 1 & 1 & 1 & 1 & 1 & 1 & 1 & 1 & 1 & 1 & 1 \\
well mixed & 1 & 30 & 20 & 1 & 1 & 1 & 1 & 1 & 1 & 1 & 1 & 1 & 1 & 1 & 1 \\
well mixed & 1 & 60 & 10 & 1 & -- & -- & 1 & 1 & 1 & 1 & 1 & -- & 1 & 1 & 1 \\
\addlinespace
well mixed & 0 & 15 & 20 & 1 & 1 & 1 & 1 & 1 & 1 & 1 & 1 & 1 & 1 & .65 & 0 \\
well mixed & 0 & 30 & 20 & 1 & 1 & 1 & 1 & 1 & 1 & 1 & 1 & 1 & 1 & 0 & 0 \\
well mixed & 0 & 60 & 10 & 1 & -- & -- & 1 & 1 & 1 & 1 & 1 & -- & .60 & 0 & 0 \\
\bottomrule

\end{tabular}
\end{table}

\begin{table}[ht]
\centering
\caption{Threshold of adult recovery $\eadult^c(\lambda,L)$ above which the colony persists: midpoint of the logistic
fit of the extinction probability by week $1\,000$ against $\eadult$. In brackets, where computed, the
profile-likelihood 95\% interval of the midpoint (width of the sigmoid profiled out; $>1$: beyond the grid).
$^\ast$Only one value of $\eadult$ gives an intermediate extinction fraction, so the width of the transition is not
identified. $^\dagger$No intermediate value: the threshold lies between the last $\eadult$ at which every colony goes
extinct and the first at which none does. $>1$: every colony goes extinct even with $\eadult=1$, or the fitted
midpoint exceeds 1. n.s.: not simulated. In space with retreat sites; well mixed without them. Sources:
\texttt{results\_em/adult\_recovery/transicion.csv} and
\texttt{results\_em/robustness\_checks/epsc\_profile.csv}.}
\label{tab:sm-epsc}
\footnotesize
\begin{adjustbox}{max width=\textwidth}
\begin{tabular}{@{}lcccc@{}}
\toprule
Structure and mating & $\lambda$ & $L=15$ & $L=30$ & $L=60$ \\
\midrule
spatial & 0 & 0.133 [0.113, 0.157] & 0.085 [0.073, 0.099] & 0.050 [0.037, 0.063]$^\ast$ \\
spatial & 0.5 & 0.22 & 0.17 & n.s. \\
spatial & 0.8 & 0.34 & 0.26 & n.s. \\
spatial & 0.9 & 0.46 [0.43, 0.50] & 0.38 [0.35, 0.42] & 0.31 [0.27, 0.35] \\
spatial & 0.95 & 0.54 & 0.48 & n.s. \\
spatial & 1 & 0.97 [0.88, $>1$] & 0.73 [0.67, 0.81] & 0.53 [0.46, 0.62] \\
\addlinespace
well mixed, cell & 0 & 0.80 [0.79, 0.86]$^\ast$ & 0.6--0.8$^\dagger$ & 0.60 [0.57, 0.67]$^\ast$ \\
well mixed, cell & 0.5 & $>1$ & $>1$ & $>1$ \\
well mixed, cell & 1 & $>1$ & $>1$ & $>1$ \\
\addlinespace
well mixed, $3\times3$ block & 0 & 0.84 & 0.60 & 0.58 \\
well mixed, $3\times3$ block & 0.5 & $>1$ & 0.91 & 0.70 \\
well mixed, $3\times3$ block & 1 & $>1$ & $>1$ & $>1$ \\
\addlinespace
well mixed, global & 0 & 0.81 [0.76, 0.86] & 0.51 [0.48, 0.54] & 0.30 [0.27, 0.33]$^\ast$ \\
well mixed, global & 0.5 & $>1$ & 0.69 & 0.40 \\
well mixed, global & 1 & $>1$ & $>1$ & $>1$ \\
\bottomrule

\end{tabular}
\end{adjustbox}
\end{table}

\paragraph{The tolerated recovery decreases with colony size}
In space the threshold decreases with $L$ for every $\lambda$, with no sign of saturation between $L=15$ and $60$:
with $\lambda=0$ it is $0.133$ [$0.113$, $0.157$], $0.085$ [$0.073$, $0.099$] and $0.050$ [$0.037$, $0.063$] at $L=15$,
30 and 60, and with $\lambda=1$ it is $0.97$ [$0.88$, $>1$], $0.73$ [$0.67$, $0.81$] and $0.53$ [$0.46$, $0.62$]
(profile-likelihood 95\% intervals of the midpoint of a binomial logistic fit; for $\lambda=0$ and $1$ the intervals
for different $L$ do not overlap; for $\lambda=0.9$ those for $L=30$ and $60$, $[0.347, 0.420]$ and $[0.273, 0.349]$,
overlap by $0.002$; likelihood-ratio $p\le0.008$ for every pair of sizes, also for $\lambda=0.9$). The width of the transition is not
identified at $L=60$, where only one or two values of $\eadult$ give an intermediate extinction fraction, and a
common width for the three sizes is not rejected ($p=0.26$, $0.11$ and $0.12$ for $\lambda=0$, $0.9$ and $1$), so the
data do not show whether the transition sharpens. Three sizes do not determine the limit of the threshold as
$L\to\infty$, which may be zero.
For any $\eadult>0$ the spatial collapse may therefore be a finite-size effect; only the edge $\eadult=0$ is
supported at every size simulated ($270$ of $270$ spatial colonies at $L=15$--$60$). Well mixed with mating in the
cell the series ($0.80$, between $0.6$ and $0.8$, and $0.60$ for $\lambda=0$) is resolved only to the grid of
$\eadult$ (likelihood-ratio $p=0.16$ and $1.0$ for the pairs of sizes); it is compatible with the value of the
deterministic limit, $0.585$, at which the deterministic orbit acquires an attracting invariant curve, but does not
show convergence to it. With global mating it decreases fast ($0.81$, $0.51$, $0.30$, non-overlapping intervals) and
that limit is close to zero. The conclusion of the main text is
conditional on the scale of the experiment for this reason: a much larger closed colony would need an even less
reversible non-breeding state.

\paragraph{A measurable indicator of the outcome}
Late breeding, the fraction of adults of breeding age that conceive from 30 weeks after the end of natality, is at
most $0.008$ in 90\% of the 819 spatial colonies that die out and at least $0.087$ in 90\% of the 621 that persist
(area under the ROC curve $0.97$; $0.99$ for $\lambda=1$; well mixed $0.99$, with a 10th percentile of $0.017$ in the
42 colonies that persist), over the sweep of $\eadult\times\lambda$ at $L=30$. Source:
\texttt{results\_em/operational\_indicators/tabla\_auc.csv}.

\paragraph{Held-out patterns of the reversible branch}
The two patterns of earlier confined colonies scored on the reversible-suppression null (main text, Section~3.2) were
defined after the production runs. The null ($\eadult=0.8$, $\lambda=0.9$, $L=30$, 200 runs) and \Cone\ were re-run
bit for bit with the archived seeds to record every litter. Each run is split into four windows: \emph{growth}, from
week 10 to the entry into the plateau; \emph{plateau}, the weeks with $N\ge0.9\,N_{\max}$ contiguous to the maximum
(median weeks 34 to about 100); \emph{decline}, from the end of the plateau to the following trough (median week
141); and \emph{late}, weeks $\ge300$. \emph{Litter loss}: a litter fails when none of its pups survives; the
fraction of failed litters is computed exactly from the litter record in each window. The weekly births per
reproductive female factor exactly as $B=C\,\bar\ell\,\varphi$ (conception rate, mean litter size at conception, a constant of the model, and
mean pup survival), so the fall of natality from growth to plateau splits additively on a logarithmic scale; the share of
conception is $\ln(C_{\rm g}/C_{\rm p})/[\ln(C_{\rm g}/C_{\rm p})+\ln(\varphi_{\rm g}/\varphi_{\rm p})]$. The null gives failed litters $0.048$ [$0.046$, $0.050$] during growth, $0.243$ [$0.239$, $0.248$] on
the plateau, $0.073$ [$0.067$, $0.079$] in the decline and $0.102$ [$0.101$, $0.102$] late (\Cone: $0.026$ and $0.087$
during growth and on the plateau); conception falls $8.4$-fold and pup survival $2.15$-fold, a conception share of
$0.73$ [$0.73$, $0.74$]. \emph{Females that never breed}: the fraction of the females born in a window that do not
conceive by 48 weeks of age, those that die earlier counting as never breeding (the definition of O11b). It is
$0.889$ [$0.885$, $0.893$] for females born on the plateau of the null ($41\,202$ females) against $0.961$ [$0.958$,
$0.964$] in \Cone, $0.366$ [$0.356$, $0.376$] for those born in the decline and $0.719$ [$0.718$, $0.720$] for those born
from week 300 (to $T-48$); in \Cone\ no female is born after the plateau. Means over runs with bootstrap 95\%
intervals. Source: \texttt{results\_em/robustness\_checks/h2h3.json}, from
\texttt{\{rev1000,c1x\}\_\{litters,cohorts\}.csv} (scripts \texttt{f10\_reversible.py} and \texttt{f10\_h2h3.py}).

\paragraph{Removal of non-breeders, background mortality and protection of juveniles}
These experiments use \Cone\ at $L=30$, in space with retreat sites and well mixed without them, with a horizon of
$1\,500$ weeks (removal) or $1\,000$ weeks (mortality). \emph{Single removal.} At week $t_r$ a fraction $f$ of the
adults with $s<s_v$ is removed at random ($f\in\{0.25,0.5,0.75,0.9,1\}$; $t_r$ from 20 to 100 weeks in space and from
15 to 50 well mixed; 20 replicates per cell). It delays extinction but does not prevent it: $12$ of $1\,000$ spatial
colonies (all with $f\ge0.9$ and $t_r=30$--$60$) and $0$ of $900$ well-mixed ones were alive at week $1\,500$; the
median extinction time rises at most from $163$ to $761$ weeks in space and from $148.5$ to $272$ well mixed. Removing
instead the adults that did not breed in the previous 8 weeks, once, gives $3$ survivors of $240$. \emph{Sustained
removal.} Removing $25\%$ or $50\%$ of the non-breeders every four weeks from $t_r=30$ keeps $20$ of $20$ spatial colonies
alive at week $1\,500$; well mixed, removing $50\%$ from $t_r=20$, 25 or 30 does the same, and $25\%$ rescues $10$--$11$
of $20$ from $t_r=20$--$25$ and none from $t_r=30$. Starting later rescues fewer ($t_r=50$ in space: $4$ and $15$ of
$20$). If the removal stops after 52, 104 or 208 weeks the colony collapses again: $2$ of $60$ spatial and $0$ of $60$
well-mixed colonies were alive at week $1\,500$. \emph{Background mortality.} An extra weekly mortality $\mu_{\rm bg}$,
independent of age and state, does not rescue \Cone: $10$ of $10$ colonies went extinct at each
$\mu_{\rm bg}\in\{0.005,0.01,0.02,0.03,0.05\}$, in space and well mixed. It rescues only when competence can
regenerate: with asocial learning ($\lambda=0$, $\eadult=0$) from $\mu_{\rm bg}=0.02$ in space ($10$ of $10$ alive) and
$0.05$ well mixed, and with $\lambda=0.9$, $\eadult=0.3$ from $0.01$ in space and $0.03$ well mixed. \emph{Protection
of juveniles.} From week $t_1\in\{15,20,25,30,40\}$, for 12, 26 or 52 weeks or permanently, the juveniles in their
window either learned from perfect demonstrators ($\bar s_V:=1$) or suffered no crowding damage (a crowding-free
nursery); 10 replicates per cell. Of the $800$ manipulated colonies $798$ went extinct; the two survivors had perfect
demonstrators permanently from week 30 or 40, in space. The median extinction time moved by at most $159$ weeks in
space and $12$ well mixed. Sources: \texttt{results\_em/hydra\_removal/summary.csv},
\texttt{results\_em/background\_mortality/summary.csv} and \texttt{results\_em/juvenile\_protection/tabla\_p3.csv}.

\paragraph{Not a mate-finding Allee effect}
In the reference model a female conceives only with a male of her own cell. Well mixed, with $\lambda=1$ and
$\eadult=0$, \Cone\ went extinct in every colony under the three mating rules, $50$ of $50$ under each
(Table~\ref{tab:sm-mate}); global mating makes the collapse earlier, not rarer. The short-lived control (SL), whose
competence regenerates individually, also dies out well mixed with mating in the cell or in the block, but with
global mating it persists for $L\ge45$ (none of $60$ colonies extinct by week $3\,000$, and no late extinction in $1.6\cdot10^5$
colony-weeks). Under perfect mate search the well-mixed variant therefore separates the two
models: the social collapse of \Cone\ does not depend on mate finding, whereas the extinction of the well-mixed SL with
local mating does.

\begin{table}[ht]
\centering
\caption{Well-mixed \Cone\ ($\lambda=1$, $\eadult=0$, no retreat sites; horizon $2\,000$ weeks) and the short-lived
control SL (horizon $3\,000$ weeks) with the father drawn in the female's cell, in her $3\times3$ block or in the
whole arena. Extinct/replicates at each $L$ (\Cone: $L=15$, $30$, $60$; SL: $L=10$, 15, 20, 30, 45, 60, 90, and 120
with global mating), the pooled count with its $95\%$ Wilson interval when every colony went extinct, the probability
of extinction by week 300 at each $L$, and the range of the median extinction time over $L$ (weeks). Source:
\texttt{results\_em/mate\_search/tabla\_mate.csv}.}
\label{tab:sm-mate}
\scriptsize
\setlength{\tabcolsep}{3pt}
\begin{adjustbox}{max width=\textwidth}
\begin{tabular}{@{}llllllc@{}}
\toprule
Model & Mating & $L$ & Extinct / replicates & Pooled & $P(\text{extinct by }300)$ & median $T_{\rm ext}$ \\
\midrule
\textsf{C1} & cell & 15--60 & 20/20, 20/20, 10/10 & 50/50 [0.93, 1.00] & 1, 1, 1 & 147--150 \\
\textsf{C1} & $3\times3$ block & 15--60 & 20/20, 20/20, 10/10 & 50/50 [0.93, 1.00] & 1, 1, 1 & 139--143.5 \\
\textsf{C1} & global & 15--60 & 20/20, 20/20, 10/10 & 50/50 [0.93, 1.00] & 1, 1, 1 & 138--141.5 \\
\addlinespace
\textsf{SL} & cell & 10--90 & 30/30, 30/30, 30/30, 30/30, 20/20, 20/20, 10/10 & 170/170 [0.98, 1.00] & 1, 1, 1, 1, 1, 1, 1 & 86.5--128.5 \\
\textsf{SL} & $3\times3$ block & 10--90 & 30/30, 30/30, 30/30, 30/30, 20/20, 20/20, 10/10 & 170/170 [0.98, 1.00] & 1, 1, 1, 1, 1, 1, 1 & 78--125.5 \\
\textsf{SL} & global & 10--120 & 30/30, 30/30, 30/30, 24/30, 0/20, 0/20, 0/10, 0/10 & -- & .93, .87, .47, .20, 0, 0, 0, 0 & 76--1070.5 \\
\bottomrule

\end{tabular}
\end{adjustbox}
\end{table}

\paragraph{Extinction is certain; its scaling distinguishes collapse from persistence}
Competence is updated before reproduction and the damage $\gamma(m_i-m_c)_+$ exceeds $1+r$ above an occupancy
$\bar m=m_c+(1+r)/\gamma$, so no animal in a cell with more than $\bar m$ occupants can breed; a pigeonhole argument
then bounds the females that can conceive in a week by a multiple of $|G|$. With no immigration, extinction of a closed
colony is certain for every variant with $\gamma>0$ (a standard argument for a population with an absorbing empty
state and a bounded number of breeders). The
fact that \Cone\ goes extinct in 200 of 200 runs is therefore informative only through the time scale. A collapse
has an extinction time that grows like $\ln|G|$; a metastable colony one that grows like $e^{c|G|}$. Both behaviours
are observed. With $\eadult=0$ the median extinction time of \Cone\ grows by $1$--$2$ weeks per unit of $\ln|G|$ well
mixed (all mating rules) and by about $6$ ($\lambda=0$) and $11$ ($\lambda=1$) weeks in space, with noisier values
between, where the last pockets die later in larger arenas (\texttt{escaleo.csv}); over the well-mixed ensembles of $L=10$--$250$ the slope is $1.25$ [$1.16$, $1.33$] weeks per
unit of $\ln N_0$, and the median residual time from the last birth to extinction is $115$--$118$ weeks at every size
(\texttt{results\_em/wellmixed\_size\_scaling/tabla\_wm.csv}). On the persistent side of the threshold the extinction probability falls with size
(space, $\lambda=1$, $\eadult=1$: $0.55$, $0.05$ and $0$ at $L=15$, $30$ and $60$), and the spatial SL is metastable,
with a late extinction rate that decreases like $e^{-c|G|}$, $c\approx0.005$ per cell. The answer to ``is the collapse
a consequence of the small size of the colony?'' is thus negative for $\eadult=0$, where the time to extinction grows
only logarithmically with size, and open for $\eadult>0$ in space, where the tolerated recovery decreases with size.

\paragraph{The mechanism cohort by cohort in the well-mixed limit}
As $|G|\to\infty$ at fixed density the well-mixed model approaches a deterministic kinetic equation for the
density of individuals by age and competence: at $L=90$ the largest difference between its density and the mean of the
simulations is $0.19\%$ of the peak (\texttt{results\_em/kinetic/resumen\_kin\_vs\_abm\_L90.csv}).
Along its orbit for \Cone\ the density $\nu$ of other individuals per cell exceeds $\nu^*=6.92$, the occupancy at which
the expected damage $\gamma\,\mathrm{E}(\mathrm{Poisson}(\nu)+1-m_c)_+$ equals the maximal gain $r$, from week 25 to
week 108. Table~\ref{tab:sm-cohorts} follows five cohorts. Their competence at the end of the window, relative to
that of their mothers, collapses between the cohorts of weeks 12 and 22, while the mean proportion $\bar\theta$ of
neighbours with $s>0$ during their window is still above the transmission threshold $\theta^*=(1-w)/(ra_w)=2/3$ for
all but the last, and while a gain bound that ignores damage would still predict ratios of $0.63$--$0.90$. The cohorts
thus fail through damage inside the window, not through a lack of demonstrators. That a sustained density above
$\nu^*$ \emph{implies} the collapse is a conjecture, verified at every point of the bifurcation diagram of the kinetic
equation but not proved. This account holds for the well-mixed model. In space the
mean crowding is far below $\nu^*$ (2.9 animals per cell at the peak), the crowding experienced by juveniles exceeds
the maximal gain only during the growth phase, and the colony collapses because the adults, damaged during the
aggregated growth and unable to recover, become non-breeding: an adult ratchet, which is why adult recovery decides
the outcome.

\begin{table}[ht]
\centering
\caption{Cohorts of the kinetic (well-mixed, deterministic) limit of \Cone. $b$: week of birth. Ratio: mean competence
of the cohort at the end of its window ($a=13$) over that of its mothers. $\bar\theta$: mean, over the window weeks
$b+1,\dots,b+12$, of the proportion of neighbours with $s>0$. Bound without damage: the ratio that the mean-field gain
alone would give. Last column: $\bar\theta\ge\theta^*=2/3$. Source:
\texttt{results\_em/kinetic/cohortes\_C1.csv}.}
\label{tab:sm-cohorts}
\footnotesize
\begin{tabular}{@{}ccccc@{}}
\toprule
$b$ & ratio & $\bar\theta$ & bound without damage & $\bar\theta\ge\theta^*$ \\
\midrule
3 & 0.991 & 1.00 & 1.02 & yes \\
12 & 0.718 & 1.00 & 0.90 & yes \\
16 & 0.188 & 0.94 & 0.77 & yes \\
19 & 0.0301 & 0.82 & 0.63 & yes \\
22 & 0.00419 & 0.63 & 0.48 & no \\
\bottomrule

\end{tabular}
\end{table}

\paragraph{Status of the claims}
The claims of the main text about the mechanism have different status. Definitions: $\nu^*=6.92$ and $\theta^*=2/3$,
computed exactly. Short argument: extinction of a closed colony is certain for $\gamma>0$ (see above; uninformative
about time scales). Conjectural: that a density above $\nu^*$ during the window implies the collapse (verified
numerically in the kinetic equation). Simulated: the level of the peak, $0.36\Kg$; growth to the numerical cap without
crowding damage (200 of 200 spatial runs); the collapse for $\eadult\le0.02$ and the threshold $\eadult^c(L)$, whose
limit $L\to\infty$ is not determined (Tables~\ref{tab:sm-pext}--\ref{tab:sm-epsc}); the effects of removing
non-breeders, of background mortality and of protecting the juveniles; the absence of a mate-finding Allee effect,
well mixed (Table~\ref{tab:sm-mate}); the logarithmic growth of the extinction time; the crossing of $\theta^*$
before the end of natality, also in colonies that persist; and the spatial mechanism, an adult ratchet during the
aggregated growth.
